\section*{Capítulo \ref{ch:g11:electric-gravitational-fields} --- Campos eléctricos y gravitatorios}

\begin{solution}{exo:g11:electric-gravitational-fields:1}
$E = \num{5.0e-3} / \num{2.0e-6} = \qty{2.5e3}{N/C}$, hacia el norte.
Sobre $q'$: $F = \num{4.0e-6} \times \num{2.5e3} = \qty{1.0e-2}{N}$,
hacia el \emph{sur} ($q' < 0$).
\end{solution}

\begin{solution}{exo:g11:electric-gravitational-fields:2}
\emph{1.} $E = \num{8.99e9} \times \num{5.0e-9} / 0.30^2
\approx \qty{5.0e2}{N/C}$, alejándose de $O$ ($Q > 0$).

\emph{2.} Para que $k\abs{Q}/d^2$ se reduzca a la mitad hace falta
duplicar $d^2$: $d' = 0.30\sqrt{2} \approx \qty{0.42}{m}$.
\end{solution}

\begin{solution}{exo:g11:electric-gravitational-fields:3}
$E = 600 / \num{3.0e-3} = \qty{2.0e5}{V/m}$;
$F = eE = \num{1.60e-19} \times \num{2.0e5} = \qty{3.2e-14}{N}$.
\end{solution}

\begin{solution}{exo:g11:electric-gravitational-fields:4}
$g_{\text{Mars}} = P/m = 930/250 = \qty{3.7}{N/kg}$. Astronauta:
$P = 80 \times 3.72 \approx \qty{3.0e2}{N}$ (unos $\qty{780}{N}$ en la
Tierra).
\end{solution}

\begin{solution}{exo:g11:electric-gravitational-fields:5}
(a) Falso: el \omterm{def:g11:electric-gravitational-fields:field}{campo} tiene una sola dirección en cada punto, así que las
líneas nunca se cortan. (b) Falso: salen del $+$ y llegan al $-$. (c)
Verdadero: el apiñamiento codifica la intensidad.
\end{solution}

\begin{solution}{exo:g11:electric-gravitational-fields:6}
\emph{1.} $d = \num{6.37e6} + \num{4.2e5} = \qty{6.79e6}{m}$, luego
$g = \num{6.67e-11} \times \num{5.97e24} / (\num{6.79e6})^2
\approx \qty{8.6}{N/kg}$ --- el $88\,\%$ del valor en la superficie.

\emph{2.} Flotan porque la estación y el astronauta caen juntos (caída
libre), no porque falte la gravedad.
\end{solution}

\begin{solution}{exo:g11:electric-gravitational-fields:7}
$E = m_e g / e = \num{9.11e-31} \times 9.81 / \num{1.60e-19}
\approx \qty{5.6e-11}{V/m}$: unas $\num{4e14}$ veces más débil que el
\omterm{def:g11:electric-gravitational-fields:field}{campo} del \omterm{def:g10:signals-and-waves:oscilloscope}{osciloscopio}. La gravedad es completamente despreciable en la
física del electrón.
\end{solution}

\begin{solution}{exo:g11:electric-gravitational-fields:8}
\emph{1.} $E = \num{5.0e3} / \num{2.0e-2} = \qty{2.5e5}{V/m}$; en
equilibrio $qE = mg$, luego
$q = \num{2.0e-6} \times 9.81 / \num{2.5e5} \approx \qty{7.8e-11}{C}$.

\emph{2.} La \omterm{def:g10:inertia:force}{fuerza} sobre la mota positiva tiene que apuntar hacia
arriba, así que $\vect E$ apunta hacia arriba: la placa positiva es la
\emph{inferior}.

\emph{3.} $n = \num{7.85e-11} / \num{1.60e-19} \approx \num{4.9e8}$
\omterm{def:g11:fundamental-interactions:elementary}{cargas elementales}.
\end{solution}

\begin{solution}{exo:g11:electric-gravitational-fields:9}
Fuera del segmento, los dos \omterm{def:g11:electric-gravitational-fields:field}{campos} apuntan en el mismo sentido; entre
las cargas se oponen --- y el equilibrio exige que el punto esté más
cerca de la fuente más débil, $+Q$. Con $x$ la distancia a $+Q$:
$kQ/x^2 = 4kQ/(0.30 - x)^2$, luego $0.30 - x = 2x$ y
$x = \qty{0.10}{m}$ desde $+Q$.
\end{solution}

\begin{solution}{exo:g11:electric-gravitational-fields:10}
$F_e = ke^2/d^2 = \num{8.99e9} \times (\num{1.60e-19})^2 /
(\num{5.3e-11})^2 \approx \qty{8.2e-8}{N}$;
$F_g = G m_p m_e / d^2 = \num{6.67e-11} \times \num{1.67e-27} \times
\num{9.11e-31} / (\num{5.3e-11})^2 \approx \qty{3.6e-47}{N}$. Cociente
$F_e / F_g \approx \num{2.3e39}$.
\end{solution}

\begin{solution}{exo:g11:electric-gravitational-fields:11}
$E = \sqrt{300^2 + 400^2} = \qty{500}{N/C}$, a
$\tan^{-1}(400/300) \approx \ang{53}$ del este hacia el norte. Sobre
$q < 0$: $F = \num{2.0e-6} \times 500 = \qty{1.0e-3}{N}$, en sentido
opuesto al \omterm{def:g11:electric-gravitational-fields:field}{campo} --- a $\ang{53}$ del oeste hacia el sur.
\end{solution}

\begin{solution}{exo:g11:electric-gravitational-fields:12}
De $GM/(R+h)^2 = \tfrac12\,GM/R^2$ sale $R + h = R\sqrt{2}$, luego
$h = (\sqrt{2} - 1)R \approx \qty{2.6e6}{m}$ (unos \qty{2600}{km}). Con
$h = R$, $d = 2R$: la cuarta parte del valor en la superficie.
\end{solution}

\begin{solution}{exo:g11:electric-gravitational-fields:13}
\emph{1.} $AM = \sqrt{3.0^2 + 4.0^2} = \qty{5.0}{cm}$; cada carga crea
$E_1 = \num{8.99e9} \times \num{2.0e-9} / 0.050^2
\approx \qty{7.2e3}{N/C}$.

\emph{2.} Las componentes paralelas a $AB$ se cancelan por simetría;
cada \omterm{def:g11:electric-gravitational-fields:field}{campo} aporta $E_1 \times 4/5$ a lo largo de la mediatriz,
alejándose del segmento:
$E = 2 \times \num{7.19e3} \times 0.80 \approx \qty{1.2e4}{N/C}$,
dirigido según la mediatriz y alejándose de $[AB]$.
\end{solution}

\begin{solution}{exo:g11:electric-gravitational-fields:14}
Las cargas libres de la caja se mueven bajo el \omterm{def:g11:electric-gravitational-fields:field}{campo} exterior: el $+$ se
acumula en una cara y el $-$ en la otra, y esas cargas desplazadas crean
un \omterm{def:g11:electric-gravitational-fields:field}{campo} interior que se opone al exterior. Siguen moviéndose hasta que
el \omterm{def:g11:electric-gravitational-fields:field}{campo} total del interior es exactamente nulo --- y entonces hay
equilibrio. La gravedad solo ofrece un signo de fuente: cada masa
\emph{añade} un \omterm{def:g11:electric-gravitational-fields:field}{campo} atractivo y ninguna puede oponerse. Un escudo
antigravedad necesitaría masa negativa, que no existe.
\end{solution}

\begin{solution}{exo:g11:electric-gravitational-fields:15}
\emph{1.} El \omterm{def:g10:universal-gravitation:weight}{peso} $m\vect g$ (hacia abajo), la tensión $\vect T$ (según
el hilo) y la \omterm{def:g10:inertia:force}{fuerza} eléctrica $q\vect E$ (horizontal). En equilibrio,
$T\sin\theta = qE$ y $T\cos\theta = mg$; dividiendo,
$\tan\theta = qE/(mg)$.

\emph{2.} $\tan\theta = \num{2.0e-7} \times \num{1.0e4} /
(\num{5.0e-4} \times 9.81) \approx 0.41$, luego
$\theta \approx \ang{22}$; $T = mg/\cos\theta \approx \qty{5.3e-3}{N}$.
\end{solution}

\begin{solution}{pb:g11:electric-gravitational-fields:1}
\textbf{1.} \omterm{def:g11:electric-gravitational-fields:capacitor}{Uniforme}, perpendicular a las placas y dirigido hacia abajo
(de la placa superior $+$ a la inferior $-$); que sea \omterm{def:g11:electric-gravitational-fields:capacitor}{uniforme} significa
que la \omterm{def:g10:inertia:force}{fuerza} sobre la gota es la misma dondequiera que derive.
\textbf{2.} $E = U/d = 800 / \num{8.0e-3} = \qty{1.0e5}{V/m}$.
\textbf{3.} La carga es negativa, así que la \omterm{def:g10:inertia:force}{fuerza} va en sentido
\emph{opuesto} a $\vect E$: hacia arriba. Solo con la placa $+$ arriba
apunta $\vect E$ hacia abajo y la \omterm{def:g10:inertia:force}{fuerza} eléctrica se pelea con el \omterm{def:g10:universal-gravitation:weight}{peso}.
\textbf{4.} $P = mg = \num{4.90e-15} \times 9.81 \approx
\qty{4.8e-14}{N}$.
\textbf{5.} $\abs{q}E = mg$.
\textbf{6.} $\abs{q} = mg/E = mgd/U$.
\textbf{7.} $\abs{q} = \num{4.81e-14} \times \num{8.0e-3} / 800
\approx \qty{4.8e-19}{C}$.
\textbf{8.} $n = \num{4.8e-19} / \num{1.60e-19} = 3$ electrones en
exceso.
\textbf{9.} $E$ crece, así que $\abs{q}E > mg$: la \omterm{def:g10:inertia:force}{fuerza} neta apunta
hacia arriba.
\textbf{10.} De $\abs{q}\,U/d = mg$ sale $U = mgd/\abs{q}$ ---
inversamente proporcional a la carga.
\textbf{11.} $\abs{q} = mgd/U = \num{3.85e-16}/1200 \approx
\qty{3.2e-19}{C} = 2e$: ha bajado de $3e$, o sea que la gota \emph{ha
perdido} un electrón.
\textbf{12.} $\abs{q} = \num{3.85e-16}/U$: \qty{3.2e-19}{C},
\qty{4.8e-19}{C}, \qty{6.4e-19}{C}, \qty{8.0e-19}{C} y
\qty{9.6e-19}{C}.
\textbf{13.} Son $2e, 3e, 4e, 5e, 6e$ con $e = \qty{1.60e-19}{C}$.
\textbf{14.} Cada escalón cambia $\abs{q}$ en exactamente
\qty{1.6e-19}{C}: los electrones llegan de uno en uno.
\textbf{15.} Haría falta $\abs{q} = \num{3.85e-16}/700 \approx
\qty{5.5e-19}{C} = 3.4\,e$ --- que no es un número entero de
electrones, así que ningún destello de radiación puede producirlo.
\textbf{16.} La carga está \emph{cuantizada}: viene en múltiplos enteros
de la \omterm{def:g11:fundamental-interactions:elementary}{carga elemental} $e = \qty{1.60e-19}{C}$ --- la magnitud que midió
Millikan.
\textbf{17.} $q \leftrightarrow m$, $E \leftrightarrow g$ y
$\abs{q}E \leftrightarrow mg$: el propio \omterm{def:g10:universal-gravitation:weight}{peso}.
\textbf{18.} $g_{\text{sphere}} = GM/d^2 = \num{6.67e-11} \times 1000 /
1.0^2 = \qty{6.7e-8}{N/kg}$; \omterm{def:g10:inertia:force}{fuerza} sobre la gota
$\num{4.90e-15} \times \num{6.7e-8} \approx \qty{3.3e-22}{N}$ --- unas
$\num{1.5e8}$ veces menor que su \omterm{def:g10:universal-gravitation:weight}{peso}. Inútil.
\textbf{19.} Como $G$ es tan pequeña, solo una masa del tamaño de un
planeta colocada encima podría rivalizar con el tirón de la Tierra ---
y ninguna masa repele. La enorme $k$ y los dos signos de la carga
permiten que una placa de sobremesa le gane al planeta en los dos
sentidos.
\textbf{20.} Un electrón: $\Delta\abs{q}/\abs{q} = e/3e \approx 33\,\%$;
una \omterm{def:g11:fundamental-interactions:ladder}{molécula} de aceite:
$\Delta m/m = \num{5e-25} / \num{4.90e-15} \approx \num{e-10}$. La
escalera de la carga mueve la tensión de equilibrio en cientos de
\omterm{def:g11:circuits-and-power:voltage}{voltios}; la escalera de la masa es invisible a cualquier tensión. Por
eso este equilibrio resuelve el \omterm{def:g11:fundamental-interactions:ladder}{átomo} de la electricidad ---
$e = \qty{1.60e-19}{C}$ --- mientras que las \omterm{def:g11:fundamental-interactions:ladder}{moléculas} del aceite se
difuminan en un continuo.
\end{solution}
